
Donc :
\[
  z=a_r.\left(\sum_{i=0}^{r-1}c_{i+1}{x'}^i+{x'}^r\right).
\]
Or :
\[
\begin{aligned}
\sum_{i=0}^{r-1}c_{i+1}{x'}^i+{x'}^r
&=(-1)^r.\left(\sum_{i=0}^{r-1}\gamma^{r-i}(E-r-1).(-1)^i{x'}^i+(-1)^r{x'}^r\right)\\
&=(-1)^r.\left(\sum_{i=0}^{r}\gamma^{r-i}(E-r-1)\gamma^i(-x')\right)\\
&=(-1)^r\gamma^r(E-r-1-x')\\
&=(-1)^r\gamma^r(F-r)\\
&=(-1)^r\lambda^r(F-1)\\
&=\lambda_{-1}(F).
\end{aligned}
\]
Si on pose $a_r=y$, on obtient donc :
\[
  z=y.\lambda_{-1}(F).
\]
Par ailleurs, $y.c_0=0$; or $c_0=(-1)^{r+1}\gamma^{r+1}(E-r-1)=\lambda_{-1}(E)$; donc $y.\lambda_{-1}(E)=0$.

\textbf{Proposition 5.11.} \uline{Avec les notations et les hypothèses de} 5.1, \uline{soit} $z\in \Gr K^{\bullet}(X)$. \uline{Pour que} $z.\bar x'=0$, \uline{il faut et il suffit qu'il existe} $y\in \Gr K^{\bullet}(S)$ \uline{tel que} :
\[
  \text{i)}\quad z=y.c^r(F);
  \qquad
  \text{ii)}\quad y.c^{r+1}(E)=0.
\]
La démonstration est identique à celle de 5.10, en se plaçant dans les gradués, et en utilisant 5.6.


\section*{6. Etude de la filtration de \texorpdfstring{$K^{\bullet}(X)$}{K^{\bullet}(X)}, $X$ ayant un faisceau ample}

Dans ce paragraphe, $X$ est un schéma ayant un faisceau ample $\underline{L}$. On rappelle que dans ce cas les $K^{\bullet}$ définis en termes de faisceaux localement libres de type fini, ou de complexes parfaits, sont égaux (IV, 2.9). On note
